\documentclass[paper=a4,fontsize=10.5pt,twoside,open=any,parskip=false,headings=standardclasses]{scrbook}
\usepackage{fontspec}
\setmainfont{IBMPlexSerif}[Extension=.otf,UprightFont=*-Regular,ItalicFont=*-Italic,BoldFont=*-SemiBold,BoldItalicFont=*-SemiBoldItalic,Scale=0.94]
\setsansfont{IBMPlexSans}[Extension=.otf,UprightFont=*-Regular,ItalicFont=*-Italic,BoldFont=*-SemiBold,BoldItalicFont=*-SemiBoldItalic,FontFace={l}{n}{*-Light},Scale=0.94]
\usepackage[english]{babel}\babelprovide[import,main]{romanian}
\usepackage{xcolor,microtype}
\definecolor{ink}{HTML}{1F2328}\definecolor{muted}{HTML}{6A737D}\definecolor{rule}{HTML}{D0D4D9}\definecolor{accent}{HTML}{963A45}
\usepackage[headsepline=false]{scrlayer-scrpage}
\setkomafont{disposition}{\sffamily\bfseries\color{ink}}
\setkomafont{pageheadfoot}{\sffamily\footnotesize\color{muted}}
\setkomafont{pagenumber}{\sffamily\footnotesize\color{muted}}
\addtokomafont{footnote}{\small}
\deffootnote{1.2em}{0pt}{\textsf{\color{accent}\thefootnotemark}\enspace}
\newcommand\ParaText{Legea din 1864 a introdus învățământul primar obligatoriu și gratuit, dar aplicarea ei a depins de resursele fiecărei comune. Rapoartele revizorilor școlari arată o realitate inegală: în unele sate școala funcționa într-o casă închiriată, în altele lipsea cu totul. Statul își asumase o misiune pe care nu o putea încă finanța.}
\newcommand\Para{\ParaText\par}
\newcommand\ParaB{Generația care a urmat a transformat școala într-un instrument al construcției naționale. Programele, manualele și formarea învățătorilor au devenit subiecte de dezbatere publică, iar Spiru Haret a legat explicit educația de dezvoltarea satului românesc.\par}
\providecommand\SN[1]{\footnote{#1}}\providecommand\MR[1]{}\providecommand\MF[1]{}
\begin{document}
\color{ink}
\Style
\author{István}\title{Școală și națiune}
\chapter{Școală, stat și națiune (1864–1918)}
\Abstract
\Para\ParaB
\section{Legea instrucțiunii publice}
\ParaText\SN{Legea instrucțiunii din 25 noiembrie 1864, art.~31. Vezi și Murgescu 2010, p.~112.}\par\ParaB
\subsection{Învățământul primar}
\MR{Comparați cu situația din Transilvania, discutată în capitolul~4.}\Para\ParaB\MF{Școală rurală din județul Vâlcea, cca.~1900.}\ParaB\Para
\section{Reformele lui Spiru Haret}
\ParaText\SN{Spiru Haret (1851–1912), ministru al instrucțiunii în trei rânduri.}\par\ParaB\Para
\subsection{Învățătorul ca agent cultural}
\Para\ParaB\Para\ParaB
\subsection{Bibliotecile populare}
\Para\ParaB\Para\ParaB\Para
\section{Bilanț la 1918}
\ParaB\Para\ParaB\Para
\end{document}
